\section{The bounded clock comparison}\label{app:clock}

We prove \eqref{eq:clock-sup} with no regularity modulus assumed
for the terminal datum. This is the local estimate that permits
the operator norm, rather than strong-operator, limit in
Proposition~\ref{prop:return-norm}.

Return to the scaled conductance $\eta=(y-m)/\sqrt a$ and the OU
expression $\mathcal A=\partial_{\eta\eta}-\eta\partial_\eta$.
The deterministic phase interval is $[0,\tau]$, with
$\tau=s(\log2)$ and scalar coefficient
$Z(s)=[m-b(x(s))]^{-1}$ from \eqref{eq:phase}.
These phase objects determine the reversed coefficient
$Z(\tau-t)$ used below.

Let $I_R=(-R,R)$, $\dd\gamma=e^{-\eta^2/2}\dd\eta$,
$E=L^2(I_R,\dd\gamma)$, and $V=H^1_0(I_R)$.
For a smooth positive coefficient $q$ with fixed upper and
lower bounds consider
\begin{equation}\label{eq:clock-general}
 q u_t=(\mathcal A-\lambda)u,\qquad
 u(t,\pm R)=0,\qquad u(0)=f .
\end{equation}
The coefficient and its needed derivatives are bounded on
the fixed rectangle; $\lambda$ ranges in a compact set.

\subsection{Construction and positive-time estimates}

Choose a smooth Dirichlet basis with finite spans dense in
$V$ and $E$, and project the initial datum in $E$. On each
finite span the $q$-weighted mass matrix is uniformly positive,
so the finite linear system has a solution. Testing with $u_N$
and using $\langle\mathcal A u,v\rangle_\gamma
=-\langle u_\eta,v_\eta\rangle_\gamma$ gives
\[
 \frac{\dd}{\dd t}\int q|u_N|^2\dd\gamma
 =-2\norm{(u_N)_\eta}_2^2
   -2\operatorname{Re}\lambda\norm{u_N}_2^2
   +\int q_t|u_N|^2\dd\gamma .
\]
Consequently
\begin{equation}\label{eq:clock-energy}
 \sup_{0\le t\le\tau}\norm{u_N(t)}_2^2
            +\int_0^\tau\norm{(u_N)_\eta}_2^2\dd t
       \le C\norm f_2^2 .
\end{equation}
Take a weak subsequence in the Hilbert space $L^2_tV$.
The $L^\infty_tE$ bound survives: testing the lower norm
bound on every measurable time set $A$ gives
$\int_A\norm u_E^2\dd t\le C\norm f_2^2|A|$.
Integrate the finite equations by parts in time against
smooth tests with values in any fixed finite span. Every
term is linear, so weak convergence passes it, including
$q_t$ and the initial datum. Density gives the full
variational equation. Strong time compactness is unnecessary.

Multiplication by $q^{\pm1}$ preserves $V$, and the equation
gives $u_t\in L^2_tV'$. Time Steklov averages justify the
energy identity: on a smaller interval the averages lie in
$H^1_tV$, so differentiating their squared $E$ norm is
legitimate. The supremum norm of the difference of two
averages is bounded by its time-average $E$ norm squared
plus twice its $L^2_tV$ norm times its derivative's
$L^2_tV'$ norm. Translation continuity makes both terms
tend to zero. Forward averages near the initial endpoint
and backward averages near the final endpoint give a
$C_tE$ representative and the limiting energy identity.
Equivalently one may apply the same argument to the
$q$-weighted norm, retaining its bounded $q_t$ term.
The initial weak datum and the energy limsup give $u(0)=f$.
The zero-data energy inequality proves uniqueness. The
entire Galerkin sequence therefore has this limit.

Testing the finite equation by $(u_N)_t$ instead gives
\[
 \frac{\dd}{\dd t}\norm{(u_N)_\eta}_2^2
       +c\norm{(u_N)_t}_2^2\le C\norm{u_N}_2^2 .
\]
Choose an earlier time with bounded gradient using
\eqref{eq:clock-energy}, then integrate. Passing weakly
gives, for each $t_*>0$,
\begin{equation}\label{eq:clock-smoothing}
 \sup_{t_*\le t\le\tau}\norm{u(t)}_V^2
       +\int_{t_*}^\tau\norm{u_t}_2^2\dd t
       \le C_{t_*}\norm f_2^2 .
\end{equation}
The equation gives $u_{\eta\eta}\in L^2$ there.
Repeated differentiated equations, with the smooth $q$
commutators, and the one-dimensional equation give the
local $C^1_tC^2_\eta$ representative on positive-time
closed subrectangles. The $V$ bound extends to each
positive time by weak $V$ compactness and strong $E$
continuity; hence the spatial output is continuous.
No compatibility at the initial Dirichlet corners is used.

For $q=1$ the equation is autonomous. Its derivative $v=u_t$
satisfies the same homogeneous equation. Choose a time
strictly before $\tau/2$ with controlled $\norm v_2$ using
\eqref{eq:clock-smoothing}, then apply the energy and
positive-time estimates to $v$. The one-dimensional
Dirichlet inequality now gives
\begin{equation}\label{eq:clock-autonomous-derivative}
 \sup_{\tau/2\le t\le\tau}
       \norm{(\mathcal A-\lambda)u_0(t)}_\infty
       \le C\norm f_2 .
\end{equation}

The maximum comparison is asserted for the weak limit,
not for pointwise values of a Galerkin approximation.
Regularize the modulus by $\sqrt{|u|^2+\delta^2}$ and test
the positive part above a scalar exponential barrier.
The extra second-derivative term has nonnegative sign
and can be dropped. The $V$ energy permits the limit
$\delta\downarrow0$. This proves the maximum bound
and its version with an inhomogeneous bounded source.

\subsection{Perturbation in supremum norm}

Take $q_\varepsilon=1+\varepsilon\eta\chi(t)$, where
$\chi(t)=Z(\tau-t)$ and
$\varepsilon R\norm\chi_\infty\le1/2$.
Put $w=u_\varepsilon-u_0$ and $h=1-q_\varepsilon$.
Its exact equation is
\[
 q_\varepsilon w_t=(\mathcal A-\lambda)w+
                          h(\mathcal A-\lambda)u_0,\qquad w(0)=0 .
\]
The forcing integration by parts reads
\[
 \langle h\mathcal A u_0,w\rangle_\gamma
   =-\int (h\overline{w_\eta}+h_\eta\overline w)
                       (u_0)_\eta\dd\gamma .
\]
Both $h,h_\eta$ are $O(\varepsilon)$ on the fixed rectangle.
Young's inequality and \eqref{eq:clock-energy} give
\begin{equation}\label{eq:clock-difference}
 \sup_{t\le\tau}\norm{w(t)}_2^2+
           \int_0^\tau\norm{w_\eta}_2^2\dd t
       \le C\varepsilon^2\norm f_2^2 .
\end{equation}
Choose $t_1\in[\tau/2,3\tau/4]$ with
$\norm{w_\eta(t_1)}_2\le C\varepsilon\norm f_2$.
The Dirichlet inequality on the bounded interval implies
$\norm{w(t_1)}_\infty\le C_R\varepsilon\norm f_2$.
On the remaining interval the divided equation has source
\[
 F=q_\varepsilon^{-1}h(\mathcal A-\lambda)u_0,\qquad
 \norm F_\infty\le C\varepsilon\norm f_2
\]
by \eqref{eq:clock-autonomous-derivative}. Modulus comparison
with $M'=C_\Lambda M+\norm F_\infty$ therefore gives
\[
 \norm{w(\tau)}_\infty
 \le e^{C_\Lambda\tau}\norm{w(t_1)}_\infty+
       \int_{t_1}^\tau e^{C_\Lambda(\tau-t)}\norm{F(t)}_\infty\dd t
 \le C\varepsilon\norm f_2
 \le C_R\varepsilon\norm f_\infty .
\]
This is \eqref{eq:clock-sup}. Rescaling the compact family
of intervals $\tau(m)$ to a unit interval makes its
constants uniform for $m\in I$.

\subsection{Identification with the stopped return}

On the bounded conductance interval the physical phase speed
is positive. Reverse the remaining phase time as in
\eqref{eq:clock-equation}. The inverse-clock conductance has
generator $q_\varepsilon^{-1}\mathcal A$ and killing rate
$\lambda/q_\varepsilon$. Apply stopped It\^o to the
constructed positive-time classical representative.
Its equation cancels the drift. The clock bounds make
physical duration at most $\tau/q_-$, and the maximum
bound controls the limiting stopped values.

There is no simultaneous terminal-phase and side exit.
To verify this, extend $q_\varepsilon$ smoothly to a
globally positive bounded coefficient. The unrestricted
inverse-clock OU law satisfies its smooth uniformly
parabolic forward equation. Interior regularity followed
by narrow continuity gives a density at each positive
phase time. A side exit exactly at the fixed terminal
phase would put this conductance at $R$ or $-R$,
an event of probability zero. Thus continuous $f$ is
attained on surviving terminal paths and the zero
Dirichlet value on earlier side exits. Dominated
convergence in the stopped identity identifies the
weak evolution with the killed return kernel.
Approximation in $E$ and the maximum bound give the
same identification for arbitrary continuous initial
data used in the comparison.

For the floor clock \eqref{eq:floor-clock}, replace $V$
by $\{w\in H^1(0,y_c):w(y_c)=0\}$ and use Lebesgue
measure. The $q$-weighted mass remains positive, and
the first energy has the favorable floor drift
$-m|w(0)|^2/2$. In the derivative estimate bound the
drift term against $w_v$ by Young's inequality before
any spatial integration by parts. The same weak
passage, uniqueness, and positive-time argument apply.
Reflection contributes $u_y\dd K=0$ in It\^o.
The corresponding nondegenerate auxiliary clock law
also excludes a tie of the absorbing upper side
and initial voltage. This completes the construction
used for the reflecting boundary in Appendix~\ref{app:boundary}.
